\chapter{Erwachsene Hunde}
\label{chap:heim:erwachsen}

\section{Verfeinerung \& Alltagsübungen}

% ── Übung: Geduld an der Tür ──────────────────────────────────────────────────
\begin{uebung}{Geduld an der Tür}{erwachsen}{Geduld}{4}
  \uebungsbeschreibung{%
    Der Hund lernt, ruhig zu warten, während die Haustür geöffnet wird. Er darf erst nach einer Freigabe nach draußen.\\[6pt]
    \textbf{Schritte:}
    \begin{enumerate}[leftmargin=*, itemsep=2pt]
      \item Hund vor der Tür in Sitz oder Platz.
      \item Tür leicht öffnen – bei Bewegung sofort schließen.
      \item Erst wenn Hund ruhig bleibt: Tür vollständig öffnen.
      \item Freigabe geben (z.\,B. „Okay!") bevor Hund nach draußen darf.
    \end{enumerate}
  }
  \uebungsvariationen{%
    \begin{itemize}[leftmargin=*, itemsep=4pt]
      \item \textbf{Klingel üben:} Übung auch bei klingelnder Türglocke üben.
      \item \textbf{Besuch:} Warten, bis Besucher ins Haus gekommen ist und Hund begrüßt hat.
      \item \textbf{Auto:} Gleiche Übung beim Aussteigen aus dem Auto.
    \end{itemize}
  }
  \uebungstrainer{%
    \begin{itemize}[leftmargin=*, itemsep=4pt]
      \item Konsequenz ist wichtiger als Schnelligkeit – lieber klein anfangen.
      \item Auf Körperspannung des Hundes achten: angespanntes Sitzen ist kein echtes Warten.
      \item Freigabekommando immer klar und einheitlich verwenden.
    \end{itemize}
  }
\end{uebung}

% ── Weitere Übungen hier einfügen ─────────────────────────────────────────────
